\honest{Explain/Worship/Ignore?}{Eliezer Yudkowsky}{2007}

A tribe is wandering the grasslands when it rains. I ask the wise man why water falls from the sky. He has never thought about it, but he answers that the sky spirits battle and their blood drips down. I ask where the sky spirits come from, and he whispers, ``From the before time.'' \nb{``Semantic Stopsigns,'' nine days earlier, opens with the same kind of dialogue: questions about where things come from, ending at a mythic origin and ``Never ask that question.'' It also traces the chain of explanations back to the Big Bang and the laws of physics.}

When you don't know why it rains, I say, you have three options. Ignore the question. Explain it. Or Worship: ``enjoy the sensation of mysteriousness.'' My example of Explain at its best is the sky spirits.

Each explanation raises the same choice again. Life is explained by chemistry, chemistry by atoms, and so on down to quarks and the Big Bang. If we go on long enough, I say, we must reach an ``Explanation That Needs No Explanation,'' and this, maybe, ``is the only explanation worth knowing.''

``There---I just hit Worship.'' I do not say which part was Worship: the claim that the chain ends somewhere, or calling that end the only explanation worth knowing. \nb{Whether the chain of explanations ends is the old question of a first cause. ``Semantic Stopsigns'' called it ``the seeming paradox of the First Cause.''}

The Explain response, I say, would have been: ``Huh, that does seem paradoxical. I wonder how the apparent paradox is resolved?'' I do not try to resolve it. \nb{``Semantic Stopsigns'' had also set it aside: ``My purpose here is not to discuss the seeming paradox of the First Cause.''}

If you find the question unimportant, or put it off until tomorrow, you have hit Ignore. ``Select your option wisely,'' I end. I give no way to tell when Ignore is the wise option.
